\documentclass[11pt]{article}

\usepackage{amsmath,amssymb}
\usepackage{xurl}
\usepackage[hidelinks]{hyperref}
\setlength{\emergencystretch}{2em}

% Shared commands for notes in docs/paper-gaps/.
% Notation follows the LDT paper (references/ldt-paper/).

\newcommand{\Lean}{\textsc{Lean}}
\newcommand{\leanid}[1]{\textsf{#1}}
\newcommand{\ghissue}[1]{\href{https://github.com/LionSR/MIPStarRE/issues/#1}{GitHub issue \##1}}
\newcommand{\ghpr}[1]{\href{https://github.com/LionSR/MIPStarRE/pull/#1}{GitHub PR \##1}}

% --- Paper-compatible notation ---
\newcommand{\F}{\mathbb{F}}
\newcommand{\N}{\mathbb{N}}
\newcommand{\C}{\mathbb{C}}
\newcommand{\tr}{\operatorname{tr}}
\newcommand{\ot}{\otimes}
\newcommand{\E}{\mathop{\bf E\/}}
\newcommand{\bx}{\mathbf{x}}
\newcommand{\by}{\mathbf{y}}
\newcommand{\bu}{\mathbf{u}}
\newcommand{\bv}{\mathbf{v}}

% Approximation relation (paper uses \approx_\eps and \simeq_\zeta)
\newcommand{\approxeps}{\approx_{\varepsilon}}
\newcommand{\simgeq}{\simeq_{\zeta}}

% Polynomial spaces (paper uses \polyfunc, \polysub, \polymeas)
\newcommand{\polyfunc}[3]{\operatorname{Poly}(#1,#2,#3)}
\newcommand{\polysub}[3]{\operatorname{PolySub}(#1,#2,#3)}
\newcommand{\polymeas}[3]{\operatorname{PolyMeas}(#1,#2,#3)}


\title{The Scalar/Tensor Bridge in the Full-Slice
\texorpdfstring{$ABAB$}{ABAB} Transport}
\author{}
\date{\today}

\newcommand{\ABAB}{\textit{ABAB}}
\newcommand{\BAB}{\textit{BAB}}
\newcommand{\bx}{\mathbf{x}}
\newcommand{\by}{\mathbf{y}}
\newcommand{\bu}{\mathbf{u}}
\newcommand{\bv}{\mathbf{v}}
\newcommand{\closenessOfIP}{\text{closenessOfIP}}

\begin{document}
\maketitle

\section{The assertion}

This note concerns the transport from the full-slice $\ABAB$ term to the
evaluated-slice $\ABAB$ term in the proof of commutativity of~$G$ in
\cite[Section~\emph{Commutativity of~$G$}]{Ji2020LowIndividualDegree}.%
\footnote{The discrepancy is tracked in
\href{https://github.com/LionSR/MIPStarRE/issues/713}{GitHub issue \#713}.
In the paper~\TeX{} source consulted here, the transport appears in
\path{references/ldt-paper/commutativity-G.tex}, lines~332--401, and
the corresponding Lean file is
\path{MIPStarRE/LDT/Commutativity/Main/Auxiliary.lean}.}
The context is the main commutativity theorem for~$G$,\footnote{Theorem
label \texttt{thm:com-main} in the cited paper.} where the
commutativity difference is expanded as
\begin{equation}\label{eq:gcomterms}
  \mathbb{E}_{\bx,\by} \sum_{g,h}
    \langle \psi |
      G^{\bx}_g G^{\by}_h G^{\bx}_g \otimes I
    | \psi \rangle
  -
  \mathbb{E}_{\bx,\by} \sum_{g,h}
    \langle \psi |
      G^{\bx}_g G^{\by}_h G^{\bx}_g G^{\by}_h \otimes I
    | \psi \rangle.
\end{equation}
The second term is called the $\ABAB$ term. The proof compares
this term in the full-slice and evaluated-slice regimes.

The paper's $\closenessOfIP$ step asserts
\begin{equation}\label{eq:gcom4}
  \mathbb{E}_{\bx,\by} \sum_{g,h}
    \langle \psi |
      G^{\bx}_g G^{\by}_h G^{\bx}_g G^{\by}_h \otimes I
    | \psi \rangle
  \approx_{\sqrt{\zeta}}
  \mathbb{E}_{\bx,\by} \sum_{g,h}
    \langle \psi |
      G^{\by}_h G^{\bx}_g G^{\by}_h \otimes G^{\bx}_g
    | \psi \rangle.
\end{equation}
The right-hand side has the form $\BAB \otimes G$ and is manifestly
positive semidefinite as a tensor product of operators. The
$\closenessOfIP$ proposition is applied with $A = G^{\bx}_g \otimes I$,
$B = I \otimes G^{\bx}_g$, and $C = G^{\by}_h G^{\bx}_g G^{\by}_h$;
the closeness between $A$ and $B$ follows from self-consistency at cost
$\zeta$ (hence $\sqrt{\zeta}$ under the $\closenessOfIP$ square root).

\section{The formal scalar endpoint}

In the formalization the top-level $\ABAB$ averages are
defined as \emph{scalar} expectations on a single
register.\footnote{The two declarations are named
\texttt{fullSliceABABAvg} and \texttt{evaluatedSliceABABAvg}, in
\path{MIPStarRE/LDT/Commutativity/Transport/FullSlice.lean}.}
Concretely, the full-slice average is
\[
  \mathbb{E}_{\bx,\by} \sum_{g,h}
    \operatorname{ev}_{\psi}\bigl(  \operatorname{leftTensor}(G^{\bx}_g G^{\by}_h G^{\bx}_g G^{\by}_h)
    \bigr),
\]
and the evaluated-slice average is
\[
  \mathbb{E}_{\bu,\bv,\bx,\by} \sum_{a,b}
    \operatorname{ev}_{\psi}\bigl(  \operatorname{leftTensor}(    G^{\bx}_{[g(\bu)=a]} G^{\by}_{[h(\bv)=b]}
        G^{\bx}_{[g(\bu)=a]} G^{\by}_{[h(\bv)=b]}
      )
    \bigr).
\]
In the evaluation $\operatorname{ev}_\psi$, a $\operatorname{leftTensor}$
is just an abbreviation for $X \otimes I$. Thus both endpoints are
scalar: they average $\operatorname{tr}(\rho_{\psi}\, (X \otimes I))$.

\section{The scalar/tensor bridge}

The paper's $\closenessOfIP$ call in~\eqref{eq:gcom4} produces a
\emph{tensor} endpoint $\BAB \otimes G$, not a scalar~$\BAB \otimes I$.
To connect the formal scalar endpoints to this tensor endpoint, the
formalization inserts a scalar-to-tensor bridge.\footnote{The bridge
lemma carries the name
\texttt{xEvaluatedFullSliceABABAvg\_to\_xEvaluatedFullSliceABABtensorAvg}
and is located in
\path{MIPStarRE/LDT/Commutativity/Main/Auxiliary.lean}.}
This bridge replaces $I$ on the right register with $G^{\bx}_g$ (or
$G^{\by}_{[h(\bv)=b]}$) and is itself an application of
$\closenessOfIP$, costing an additional~$\sqrt{\zeta}$.

Write $A^{\ABAB}_{\text{full}}$ and $A^{\ABAB}_{\text{eval}}$ for the
two scalar averages displayed above. The full transport
lemma\footnote{The corresponding declaration is
\texttt{fullSlice\_scalar\_marginalize\_y} in the same file.} then
assembles the chain:
\begin{align*}
  \bigl|
    A^{\ABAB}_{\text{full}}
    - A^{\ABAB}_{\text{eval}}
  \bigr|
  &\leq
  \underbrace{\frac{md}{q} + \sqrt{\zeta}}_{\text{paper prefix}}
  +
  \underbrace{2\sqrt{\zeta}}_{\text{scalar/tensor bridge}}
  +
  \underbrace{\frac{md}{q} + \sqrt{\zeta}}_{\text{paper tail}} \\
  &= 2\cdot\frac{md}{q} + 4\sqrt{\zeta}.
\end{align*}
Without the scalar/tensor bridge, the bound would be
$2md/q + 2\sqrt{\zeta}$, matching the paper's implicit transport cost
before the commutativity-bridge step. The extra $2\sqrt{\zeta}$
reflects the price of applying $\closenessOfIP$ twice more to shuttle
between scalar and tensor endpoints.

\section{Impact on the final estimate}

The factor $2\sqrt{\zeta}$ is absorbed into the existing error budget.
The paper's total transport error without scalar/tensor overhead is
$2md/q + 2\sqrt{\zeta}$ for the $\ABAB$ term (one-way, before the
factor~$2$ from~\eqref{eq:gcomterms}). The augmented error
$2md/q + 4\sqrt{\zeta}$ is still bounded by $2md/q + 2\sqrt{\zeta}$
when $\zeta$ is replaced by a $4$-times-larger parameter, which does
not change the asymptotic exponent in the final soundness cascade.

In the public interface of the formalization, all four scalar
averages are exposed with the augmented bounds, while the tensor-form
intermediates are kept private, consistent with the decision to keep
the public interface purely scalar.\footnote{The public scalar
averages are \texttt{fullSliceABAAvg}, \texttt{evaluatedSliceABAAvg},
\texttt{fullSliceABABAvg}, and \texttt{evaluatedSliceABABAvg}; the
tensor intermediates are \texttt{xEvaluatedSliceBABAtensorAvg} and
\texttt{xEvaluatedFullSliceABABtensorAvg}.}

\section{Conclusion}

\noindent\textbf{Verdict:} a genuine quantitative overhead, but no
mathematical discrepancy. The formalization pays an extra
$2\sqrt{\zeta}$ to cross between the scalar public interface and the
tensor form required by $\closenessOfIP$. This overhead is accounted for in
the stated error bounds and does not obstruct the final soundness
estimate. The divergence is present in every lemma that transports a
scalar $\ABAB$ average through a tensor $\closenessOfIP$ step, and it
is a consequence of a deliberate design choice to keep the public API
in scalar form.

\bibliographystyle{alpha}
\bibliography{docs/paper-gaps/references}

\end{document}